% =============================================================================
% PARTE I: La Mente Inaspettata
% =============================================================================

\addcontentsline{toc}{part}{PARTE I: La Mente Inaspettata} % Questo la aggiunge all'indice

\cleardoublepage
\thispagestyle{empty}

\vspace*{\fill}

\begin{center}
    % Numero della parte
    {\sffamily\huge\textls[100]{\MakeUppercase{Parte I}}}
    
    \vspace{1.5em}
    
    % Linea decorativa
    \rule{0.5\textwidth}{0.4pt}
    
    \vspace{1.5em}
    
    % Titolo della parte
    {\sffamily\Huge\bfseries La Mente Inaspettata}
    
    \vspace{2em}
    
    % Sottotitolo/descrizione
    \begin{minipage}{0.7\textwidth}
        \centering
        {\large\itshape\color{moderngray}
        Viaggio alle origini dei nostri pensieri.\\
        Capire l'architettura mentale con cui siamo nati.}
    \end{minipage}
    
    \vspace{1.5em}
    
    % Linea decorativa
    \rule{0.5\textwidth}{0.4pt}
\end{center}

\vspace*{\fill}

\clearpage
\thispagestyle{empty}
\vspace*{2em}
\begin{center}
	{\sffamily\large\bfseries\color{darkgray} Da dove partiamo}
\end{center}
\vspace{1em}

Il prologo ci ha lasciato con una sala controllo e una domanda. Perché è così difficile vedere ciò che contraddice le nostre convinzioni, anche quando ne va di mezzo tutto? Prima di chiederci che cosa succede quando le macchine imparano da noi, dobbiamo capire che cosa c'è da imparare, cioè come una mente umana costruisce la propria versione del mondo.

Questa parte segue quella costruzione dal basso. Si comincia dall'architettura del cervello, messa insieme in centinaia di milioni di anni senza un progetto, e si passa subito alla memoria, perché è lì che il passato diventa identità e l'identità diventa convinzione. Poi vedremo i due modi in cui la mente lavora, uno rapido e uno lento, e scopriremo che il secondo spesso giustifica ciò che il primo ha già deciso. Vedremo che cosa accade quando una prova contraddice chi crediamo di essere, e infine come le stesse tendenze si moltiplicano, o si correggono, quando pensiamo in gruppo.

Alla fine di questa parte non avremo un elenco di errori. Avremo un'idea più precisa di un processo, quello per cui ricordo, emozione e appartenenza producono una realtà su misura. È il materiale grezzo di tutto ciò che segue. Le decisioni della Parte II, gli archivi da cui impareranno le macchine della Parte III, le difese della Parte IV partono tutti da qui.

\vspace{1em}
\begin{storia}
Giulia fa il medico di famiglia da vent'anni. Il primo anno, una sera d'inverno, ha mandato a casa un ragazzo con un dolore al petto. Le era sembrato agitato, giovane, sano. «Stress», aveva scritto. Il ragazzo è tornato due giorni dopo con un infarto ed è sopravvissuto. Giulia non ne parla quasi mai, ma quella sera è diventata una frase che si ripete da sola, \textit{io non sbaglio due volte}. Da allora, davanti a un uomo giovane con un dolore toracico, è la più veloce del distretto. È un buon medico. E ha un ricordo che decide al posto suo più spesso di quanto immagini.
\end{storia}

\cleardoublepage
